\scanpage{91}
% Typescript inverted in the scan.
\mspage{\handwritten{43}}
supposer $\mathrm{X}$ affine et qu'il existe un préschéma affine $\mathrm{X'_{\handwritten{1}}}$, de
présentation finie, quasi fini et fidèlement plat sur $\mathrm{X}$, et un
$\mathrm{X}$-morphisme $\mathrm{X'_{\handwritten{1}} \handwritten{\longrightarrow} X'}$, ce qui nous ramène au cas où $\mathrm{f}$ est quasi\handwritten{-}fini.
Quitte à localiser encore sur $\mathrm{X}$, au sens de la topologie étale
cette fois, on voit qu'il existera un ouvert $\mathrm{X'_{\handwritten{1}}}$ de $\mathrm{X'}$ tel que
$\mathrm{X'_{\handwritten{1}} \handwritten{\longrightarrow} X}$ soit fini et surjectif, ce qui nous ramène au cas où $\mathrm{f}$ est
fini et surjectif, déjà traité dans a).

9.4.4. Dans le cas d), on peut supposer (en se localisant sur $\mathrm{X}$)
que $\mathrm{X}$ est réduit à un seul point, et même contenu de 1.1. q\deleted{\unsure{g|y}}u'il est
spectre d'un corps $\mathrm{k}$, puis (quitte à remplacer $\mathrm{X'}$ par un ouvert
affine non vide) que $\mathrm{X'}$ est affine et à fortiori $\mathrm{f}$ est quasi compact
et quasi séparé. Mais alors il est encore vrai que $\mathrm{f_{\unsure{*|x}}}$ commute
\margin{\handwritten{$\mathrm{\uparrow}$}}
à tout changement de base $\mathrm{Y \handwritten{\longrightarrow} X}$, comme nous verrons ultérieurement.
De plus, $\mathrm{f}$ est universellement ouvert (EGA IV 2.4.9.) donc on
est s\handwritten{ous} les conditions de 9.1. , et le raisonnement de 9.4.1. s'applique
encore.

\handwritten{Ceci achève la démonstration de 9.4.}
% archive sheet number: 90

